\section{Communication au patient}

\subsection{Endométriome}
Endométriome est une maladie complexe. Il s’agit d’un type d’endométriose. Cela signifie qu’au niveau des ovaires, nous retrouvons un kyste rempli de liquide. Ce kyste contient entre autres de l’endomètre qui est la muqueuse de l’utérus, ce qui est une localisation anormale. La douleur qui revient chaque mois s’explique parce que le kyste réagit comme l’utérus au cours du cycle menstruel, c’est-à-dire qu’il se construit, prolifère au cours de la première phase du cycle et s’auto-détruit à la fin du cycle ce qui correspond aux règles au niveau de l’utérus. Cependant, le tissu utérin étant au niveau du kyste ne peut pas s’évacuer, le sang va donc s’accumuler et provoquer des douleurs intenses. L’endométriome est une lésion bégnine, mais pouvant entraîner des complications fonctionnelles importantes, notamment sur la fertilité.

\subsection{Symptômes}

%\begin{singlespace}
Les symptômes typiques sont:

\begin{itemize}[label=\textbullet]
    \item Douleurs intenses lors des menstruations qui impact le quotidien de la patiente;
    \item Douleurs pelviennes;
    \item Douleurs lors des rapports sexuels;
    \item Douleurs lors de l’évacuation des selles ou parfois des urines;
    \item Troubles de la fertilité;
\end{itemize}
%\end{singlespace}

La présence de ces symptômes ne signifie pas systématiquement la présence d’endométriose mais justifie une évaluation avec un gynécologue expérimenté. Il convient de noter que tous ces symptômes ne sont pas toujours présents et qu’ils fluctuent énormément selon chaque patiente. Dans ce cas, la patiente ressent de vives douleurs lors de ses menstruations depuis de longtemps et qui se sont accentuées ces 4 derniers mois.

\subsection{Diagnostic}
Le diagnostic de cette patiente a pu être posé en observant l’évolution du kyste appréciable grâce aux différentes imageries réalisées et à la prise de la pilule Dimetrum\textregistered.\, De plus, le marqueur CA 125, étant élevé, confirme notre diagnostic.

\subsection{Traitement}
Le traitement consiste en une intervention au cours de laquelle le chirurgien enlève toutes les lésions endométriales visibles. L’opération ne laisse que trois petits trous au niveau abdominal utiles pour le passage des instruments. Ensuite, le ventre de la patiente va être gonflé d’air afin d’avoir une meilleure vision. Une fois le kyste visible, le chirurgien pourra commencer sa résection. Si d’autres lésions sont visibles, le chirurgien en profitera pour faire un traitement global. Si l’endométriome est trop conséquent et que les ovaires risquent d’être trop endommagés, la chirurgie se fera en deux temps. Entre ces deux interventions, la patiente devra prendre un traitement hormonal afin de réduire la taille des lésions restantes ce qui facilitera la seconde chirurgie.

\subsection{Suivi}
Le suivi commence par une consultation post-opératoire afin de vérifier les plaintes éventuelles de la patiente ainsi que l’évolution de la cicatrisation au niveau ovarien. Une imagerie de contrôle sera réalisée entre 3 à 6 mois ainsi qu’une réévaluation de la réserve ovarienne à posteriori. D’autres suivis peuvent être organisés selon la volonté de la patiente tel qu’un suivi de fertilité, un suivi hormonal s'il y a un désir d’aménorrhée.